%% kapa: exam=1 round=30 year=2019 period=1 form=A subject=civil_law subject_ko=민법 seq=12 number=12
%% source: Q-Net 감정평가사 기출문제 (cst00309) artlSeq=5205205 file=extracted/1008/5205205/2019년 제30회 감정평가사 1차시험 문제지/2019년 제30회 감정평가사 1차시험 문제지/2019년도 제30회 감정평가사 1차 1교시 A형.hwp
%% conversion: kordoc 4.24.51 -> tools/gen1.py (auto); review=proofread-sonnet; flags=none
\begin{question}{12}
甲은 자신의 점포를 32만 달러에 팔기로 의욕하였지만, 미국인 乙에게 실수로 매매대금을 23만 달러로 표시하여 이 가격으로 계약이 체결되었다. 이 사안에 관한 설명으로 옳은 것은?
\begin{choices}
  \item 위 매매계약은 甲의 진의 아닌 의사표시로서 일단 유효하지만, 甲이 乙의 악의 또는 과실을 입증하여 무효를 주장할 수 있다.
  \item 甲과 乙은 모두 통정허위표시에 따른 무효를 주장할 수 있다.
  \item 甲은 오표시무해의 원칙을 주장하여 ‘32만 달러’를 대금으로 하는 매매계약의 성립을 주장할 수 있다.
  \item 甲은 착오를 주장하여 위 매매계약을 취소할 수 있지만, 乙이 甲의 중대한 과실을 증명하면 취소할 수 없다.
  \item 위 매매계약은 불합의에 해당하므로, 매매계약 자체가 성립하지 않는다.
\end{choices}
\end{question}
\answer{4}
